\documentclass[10pt,a4paper]{amsart}
\usepackage[margin=19mm]{geometry}
\usepackage{amsmath,amssymb,mathtools}
\usepackage[scr=rsfs,cal=euler]{mathalfa}
\usepackage{xfrac}
\usepackage[unicode,hidelinks,bookmarks=false]{hyperref}
\newcommand{\ie}{i.e.\ }
\newcommand{\cf}{cf.\ }
\newcommand{\resp}{respectively\ }
\newcommand{\Subsection}{\subsection}
\providecommand{\nobreakdash}{\nobreak\hbox{-}}
\setlength{\emergencystretch}{2em}
\multlinegap0pt
\usepackage{mathtools,tikz}
\usepackage{bbm}
\usepackage{enumitem}
\usepackage{tikz-cd}
\usepackage{cleveref}
\usepackage{thm-restate}
\usepackage{stmaryrd}
\newcommand{\N}{\mathbb N}
\newcommand{\cube}[1]{\{0,1\}^{#1}}
\newcommand{\bits}{\{0,1\}}
\newcommand{\unif}[1]{\mathsf H_{#1}}
\newcommand{\Pow}{\mathcal P}
\newcommand{\eps}{\varepsilon}
\newcommand{\ind}{\mathbf 1}
\newcommand{\suchthat}{\;:\;}
\newcommand{\VCdim}{\operatorname{VCdim}}
\newcommand{\copyof}[1]{\text{a copy of }#1}
\DeclareMathOperator{\diam}{diam}
\newtheorem{theorem}{Theorem}
\newtheorem{conjecture}[theorem]{Conjecture}
\newtheorem{problem}[theorem]{Problem}
\newtheorem{lemma}[theorem]{Lemma}
\newtheorem{corollary}[theorem]{Corollary}
\newtheorem{proposition}[theorem]{Proposition}
\newtheorem{claim}[theorem]{Claim}
\newtheorem{definition}[theorem]{Definition}
\newtheorem{remark}[theorem]{Remark}
\newtheorem{observation}[theorem]{Observation}
\begin{document}
\title[Ramsey Obstructions to Disambiguation]{Ramsey Obstructions to Disambiguation}
\author{Romain Bourneuf, Antonin Kiladjian, Stéphan Thomassé}
\date{}
\maketitle
\noindent\emph{Source and licence.} Ramsey Obstructions to Disambiguation. Version arXiv:2609.16359v1 (CC BY 4.0). This material is distributed under the Creative Commons Attribution 4.0 International licence, http://creativecommons.org/licenses/by/4.0/, as recorded in the arXiv OAI licence metadata for this exact version and shown on the abstract page https://arxiv.org/abs/2609.16359v1. Full rights notice: licenses/arxiv-2609.16359-CC-BY-4.0.txt.\par\smallskip
\noindent\textit{Editorial scope.} Counted unit: the original \texttt{theorem} environment \texttt{lem:palvolgyi} at source lines 623--625 together with its complete proof at lines 636--673; the delivered document keeps source lines 463--674 so that every referenced label and every citation resolves inside it. Native editable source is the paper's own concatenated TeX; the AMS wrapper is editorial formatting only. No publisher class, upstream script or unrelated artwork is redistributed.
\par\smallskip\noindent\textit{Reference note.} This article ships no compiled bibliography (biblatex); the entries delivered here are composed from the author's own BibTeX (.bib) fields, with no field invented and no entry dropped.
\section{Metric obstructions to disambiguation}

We follow the approach described in the proof overview.
We first show that a coloring satisfying a gap property on a suitable discrete torus has large VC-dimension whenever its color depends only on the coordinatewise distances. 
We then use Pálvölgyi’s Dense Block theorem to obtain such a coloring from an arbitrary disambiguation. 
An isometric embedding of even cycles into hypercubes completes the proof of the Gap-Hamming Disambiguation theorem. 
We conclude with the spherical and toroidal consequences.

\subsection{Shattering from distance profiles}

Let $(X, d)$ be a metric space. 
The \emph{diameter} of $(X, d)$ is the supremum of $d(x, y)$ over all $x, y \in X$. We denote it by $\diam_d(X)$, or simply $\diam(X)$ when $d$ is clear from context.
For $0 < \varepsilon < 1/2$, a coloring $c : X \times X \to \{0, 1\}$ is an \emph{$\varepsilon$-gap coloring} if $c(x, y) = 0$ whenever $d(x, y) \leq \varepsilon \cdot \diam(X)$ and $c(x, y) = 1$ whenever $d(x, y) \geq (1-\varepsilon) \cdot \diam(X)$.

Let $Q_n$ denote the $n$-dimensional hypercube $\{0, 1\}^n$, and let $d_{Q_n}$ denote the Hamming distance over $Q_n$.
For $S \subseteq [n]$, we denote by $\mathbbm{1}_S$ the vector $x \in Q_n$ such that $x_i = 1$ if and only if $i \in S$. 
Let $\mathbb{Z}_r = \mathbb{Z}/r\mathbb{Z}$ and for $x, y \in \mathbb{Z}_r$, let $d_{\mathbb{Z}_r}(x, y)$ be the length of the shorter arc between them on the cycle $\mathbb Z_r$. 
Formally, $d_{\mathbb{Z}_r}(x, y) = \min\{|x-y|, r - |x-y|\}$, using representatives in $\llbracket0, r-1\rrbracket$.

Let $(\mathbb Z_r)^n$ denote the $n$-dimensional discrete torus of side length $r$.
For $x, y \in (\mathbb{Z}_r)^n$, the distance between $x$ and $y$ is $d_{(\mathbb{Z}_r)^n}(x, y) = \sum_{i \in [n]}d_{\mathbb{Z}_r}(x_i, y_i)$.
Observe that $\diam((\mathbb{Z}_r)^n) = n \cdot \diam(\mathbb{Z}_r) = n \lfloor r/2 \rfloor$.
For $x, y \in (\mathbb{Z}_r)^n$, we call the vector $\delta(x, y) = (d_{\mathbb{Z}_r}(x_1, y_1), \ldots, d_{\mathbb{Z}_r}(x_n, y_n))$ the \emph{distance profile} of $(x, y)$.

A $2$-coloring $c : (\mathbb{Z}_{r})^n \times (\mathbb{Z}_{r})^n \to \{0, 1\}$ is \emph{distance-profile-canonical} if there exists a function $f : \{0, \ldots, \lfloor r/2 \rfloor\}^n \to \{0, 1\}$ such that for all $x,y \in (\mathbb{Z}_{r})^n$, we have $c(x, y) = f(\delta(x, y))$. In other words, the color of $(x, y)$ only depends on the distance profile of $(x, y)$.

Given a set $V$ and a $2$-coloring $c : V \times V \to \{0, 1\}$, we say that a set $U \subseteq V$ is \emph{shattered} by $c$ if for every subset $W \subseteq U$, there exists $w \in V$ such that for every $u \in U$ we have $c(w, u) = 1$ if and only if $u \in W$.
The \emph{VC-dimension} of $c$ is the supremum of the sizes of its finite shattered sets.
Note that the VC-dimension of $c$ is the VC-dimension of the binary matrix corresponding to $c$.

We first prove that any $\varepsilon$-gap distance-profile-canonical $2$-coloring of a discrete torus of large enough side length and dimension has large VC-dimension.
Intuitively, we choose a $0$-colored profile that is as far as possible from the origin in the torus. Every coordinate that has not yet reached the maximum then gives a direction along which increasing that coordinate changes the color from $0$ to $1$; the gap condition guarantees that there are many such directions. 
The key point is that these many independent color-changing directions can be used to encode arbitrary subsets, and hence to produce a large shattered set.

\begin{lemma}[Shattering from distance profiles]\label{lem:large-vc}
    For all $0 < \varepsilon < 1/2$ and $q \in \mathbb{N}$, there exist an integer $n$ and an even integer $R$ such that every $\varepsilon$-gap distance-profile-canonical coloring $c : (\mathbb{Z}_{R})^n \times (\mathbb{Z}_{R})^n \to \{0, 1\}$ has VC-dimension at least $q$.
\end{lemma}

\begin{proof}
    Let $s$ be large enough so that $\frac{1}{3^s} \leq \varepsilon$.
    Let $r = 3^s$ and $R = 2r$.
    Let $n$ be large enough so that $\varepsilon n \geq q$.
    The factor $3$ comes from a simple distance gadget: $a$ and $3a$ are equally far from $2a$, but have different distances from $0$.
    
    Let $c : (\mathbb{Z}_{R})^n \times (\mathbb{Z}_{R})^n \to \{0, 1\}$ be an $\varepsilon$-gap distance-profile-canonical coloring.
    Since $c$ is distance-profile-canonical, there exists a function $f : \{0, \ldots, r\}^n \to \{0, 1\}$ such that for all $x, y \in (\mathbb{Z}_{R})^n$, we have $c(x, y) = f(\delta(x, y))$.
    
    For every $i \in \{0, \ldots, s\}$, let $d_i = 3^i$; in particular $d_0 = 1$ and $d_s = r$.
    Let $x, y \in (\mathbb{Z}_{R})^n$.
    If $\delta(x, y) = (d_s, \ldots, d_s)$ then $d(x, y) = rn = \diam((\mathbb{Z}_{R})^n)$.
    Since $c$ is an $\varepsilon$-gap coloring, we then have $c(x, y) = 1$, so $f(d_s, \ldots, d_s) = 1$.
    If $\delta(x, y) = (d_0, \ldots, d_0)$ then $d(x, y) = n \leq \varepsilon rn = \varepsilon \cdot  \diam((\mathbb{Z}_{R})^n)$.
    Since $c$ is an $\varepsilon$-gap coloring, we then have $c(x, y) = 0$, so $f(d_0, \ldots, d_0) = 0$.
    
    Consider $a = (a_1, \ldots, a_n) \in \{1, 3, \ldots, r\}^n$ such that $f(a) = 0$ and $\sum_{j = 1}^n a_j$ is maximum with this property.
    Let $J = \{j \in [n] : a_j < r\}$.
    We show that $|J| \geq q$.
    If $x = (0, \ldots, 0)$ then $\delta(x, a) = a$ so \[d(x, a) = \sum_{j = 1}^n a_j \geq \sum_{j \notin J}r + \sum_{j \in J}1 = (n-|J|) \cdot r + |J| = rn - |J| \cdot (r-1).\]
    However, we have $0 = f(a) = f(\delta(x, a)) = c(x, a)$. 
    Since $c$ is an $\varepsilon$-gap coloring, this implies $d(x, a) \leq (1-\varepsilon) \cdot \diam((\mathbb{Z}_{R})^n) = (1-\varepsilon)rn$.
    Thus, $rn - |J| \cdot (r-1) \leq (1-\varepsilon)rn$ so $\varepsilon rn \leq (r-1) |J| \leq r|J|$ so $|J| \geq \varepsilon n \geq q$.
    
    For every $i \in J$, let $x^i \in (\mathbb{Z}_{R})^n$ be the vector such that \[x^i_j = \begin{cases}
      2a_j, &j \in J \setminus \{i\} \\
      0, & \text{ otherwise}.
    \end{cases}\]
    These points are pairwise distinct by definition.
    For every set $S \subseteq J$, let $y^S\in (\mathbb{Z}_{R})^n$ be the vector such that \[y^S_j = \begin{cases}
      3a_j, &j \in S \\
      a_j, &j \notin S.
    \end{cases}\]
    
    Note that for each $j \in J$, we have $a_j \leq d_{s-1} = r/3$ so $0, a_j, 2a_j, 3a_j$ all lie on an arc of length at most $r$ in $\mathbb{Z}_R$. Consequently, $d(0, a_j) = d(a_j, 2a_j) = d(2a_j, 3a_j) = a_j$ and $d(0, 3a_j) = 3a_j$.
    
    We now argue that for every $S \subseteq J$ and every $i \in J$, we have \[\delta(y^S, x^i) = \begin{cases}
        a+2a_ie_i, &i \in S \\
        a, &i \notin S.
    \end{cases}\] 
    By maximality of $a$, this will imply that $c(y^S, x^i) = 1$ if and only if $i \in S$.
    Let $S \subseteq J$ and let $i \in J$.
    Let $j \in [n]$. We consider several cases. \begin{itemize}
        \item If $j \notin J$ then $d(y^S_j, x^i_j) = d(a_j, 0) = d(r, 0) = r = a_j$.
        \item If $j \in J \setminus \{i\}$ then $y^S_j \in \{a_j, 3a_j\}$ and $x^i_j = 2a_j$ so $d(y^S_j, x^i_j) = a_j$.
        \item If $j = i$ and $i \in S$ then $y^S_j = 3a_j$ and $x^i_j = 0$ so $d(y^S_j, x^i_j) = 3a_j$.
        \item If $j=i$ and $i \notin S$ then $y^S_j = a_j$ and $x^i_j = 0$ so $d(y^S_j, x^i_j) = a_j$.
    \end{itemize}
    
    Thus, if $i \in S$, we have $\delta(y^S, x^i) = (a_1, \ldots, a_{i-1}, 3a_i, a_{i+1}, \ldots, a_n)$, so by maximality of $(a_1, \ldots, a_n)$, we have $c(y^S, x^i) = f(a_1, \ldots, a_{i-1}, 3a_i, a_{i+1}, \ldots, a_n) = 1$.
    However, if $i \notin S$, we have $\delta(y^S, x^i) = (a_1, \ldots, a_n) = a$ so $c(y^S, x^i) = f(a) = 0$.
    
    Therefore, the set $\{x^i : i \in J\}$ is shattered by $c$, so $c$ has VC-dimension at least $q$.
\end{proof}

The following two lemmas describe how gap constraints and VC-dimension behave when we pull back the two arguments of a coloring along possibly different maps.

\begin{lemma}[Gap preservation under pullback]\label{lem:transfer-gap-coloring}
    Let $0 < \varepsilon < 1/2$ and let $(X, d_X), (Y, d_Y)$ be two finite metric spaces such that there exist two maps $\iota_1, \iota_2 : Y \to X$ that satisfy $d_X(\iota_1(y), \iota_2(y')) = \lambda \cdot d_Y(y, y') + \Delta$ for all $y, y' \in Y$, for some $\lambda, \Delta \geq 0$.
    Let $c : X \times X \to \{0, 1\}$ be an $\varepsilon$-gap coloring and define $\widehat{c} : Y \times Y \to \{0, 1\}, (y, y') \mapsto c(\iota_1(y), \iota_2(y'))$.
    Suppose that $\lambda \cdot \diam(Y) \geq (1-\varepsilon/2) \cdot \diam(X)$.
    Then, $\widehat{c}$ is an $\varepsilon/2$-gap coloring.
\end{lemma}

\begin{proof}
    Let $y, y' \in Y$ such that $d_Y(y, y') = \diam(Y)$.
    Then, \[d_X(\iota_1(y), \iota_2(y')) = \lambda \cdot d_Y(y, y') + \Delta \geq \lambda \cdot \diam(Y),\] so $\diam(X) \geq \lambda \cdot \diam(Y)$.
    Furthermore, we have \[\diam(X) \geq d_X(\iota_1(y), \iota_2(y')) = \lambda \cdot d_Y(y, y') + \Delta = \lambda \cdot \diam(Y) + \Delta \geq (1-\varepsilon/2) \cdot \diam(X) + \Delta,\] so $\Delta \leq \varepsilon/2 \cdot \diam(X)$.

    Let $y, y' \in Y$ such that $d_Y(y, y') \leq \varepsilon/2 \cdot \diam(Y)$.
    Then, \begin{align*}
        d_X(\iota_1(y), \iota_2(y')) &= \lambda \cdot d_Y(y, y') + \Delta \leq \lambda \cdot \varepsilon/2 \cdot \diam(Y) + \varepsilon/2 \cdot \diam(X) \\
        &\leq \varepsilon/2 \cdot \diam(X) + \varepsilon/2 \cdot \diam(X) = \varepsilon \cdot \diam(X).
    \end{align*} Thus, since $c$ is an $\varepsilon$-gap coloring, we have $\widehat{c}(y, y') = c(\iota_1(y), \iota_2(y')) = 0$.
    \smallskip
    
    Let $y, y' \in Y$ such that $d_Y(y, y') \geq (1-\varepsilon/2) \cdot \diam(Y)$.
    Then, \begin{align*}
        d_X(\iota_1(y), \iota_2(y')) &= \lambda \cdot d_Y(y, y') + \Delta \geq \lambda \cdot (1-\varepsilon/2) \cdot \diam(Y) \\
        &\geq (1-\varepsilon/2) \cdot (1-\varepsilon/2) \cdot \diam(X) \geq (1 - \varepsilon) \cdot \diam(X).
    \end{align*} 
    Thus, since $c$ is an $\varepsilon$-gap coloring, we have $\widehat{c}(y, y') = c(\iota_1(y), \iota_2(y')) = 1$.
    This concludes the proof that $\widehat{c}$ is an $\varepsilon/2$-gap coloring.
\end{proof}

\begin{lemma}[VC-dimension preservation under pullback]\label{lem:transfer-VC-dim}
    Let $X, Y$ be two sets and let $\iota_1, \iota_2 : Y \to X$ be two maps.
    Let $c : X \times X \to \{0, 1\}$ be a coloring and define $\widehat{c} : Y \times Y \to \{0, 1\}, (y, y') \mapsto c(\iota_1(y), \iota_2(y'))$.
    If $\widehat{c}$ has VC-dimension at least $q \in \mathbb{N}$ then so does $c$.
\end{lemma}

\begin{proof}
    Suppose that $\widehat{c}$ has VC-dimension at least $q$.
    Let $\{y^i : i \in [q]\} \subseteq Y$ be shattered by $\widehat{c}$: for every $S \subseteq [q]$, there exists $y^S \in Y$ such that for every $i \in [q]$, we have $\widehat{c}(y^S, y^i) = 1$ if and only if $i \in S$.
    By definition of $\widehat{c}$, this means that for every $S \subseteq [q]$ and every $i \in [q]$, we have $c(\iota_1(y^S), \iota_2(y^i)) = 1$ if and only if $i \in S$.
    Thus, the set $\{\iota_2(y^i) : i \in [q]\}$ is shattered by $c$.
    In particular, this implies that its elements are pairwise distinct, so $c$ has VC-dimension at least $q$.
\end{proof}

\subsection{Canonicalization on a torus}

In view of \cref{lem:large-vc}, we would like to find a restriction on which the color depends only on the orbit of each coordinate pair. 
We also need the variable coordinates to occupy almost the entire ambient space, so that the gap condition survives the restriction.
Our key tool to find this is a result of Pálvölgyi \cite{Palvolgyi26}. First, we need to introduce some notation.

Let $G$ be a finite group acting on a finite set $X$.
A \emph{$G$-block map} of \emph{dimension} $m$, \emph{length} $n$, and \emph{block size} $k$ is a map $\phi : X^m \to X^n$ for which there are pairwise disjoint sets $I_1, \ldots, I_m \subseteq [n]$, each of cardinality $k$, labels $\gamma_{\ell} \in G$ for $\ell \in I_1 \cup \ldots \cup I_m$, and fixed letters $z_\ell \in X$ outside this union, such that for every $x \in X^m$, we have \[
  \phi(x)_\ell=
  \begin{cases}
    \gamma_\ell x_j,&\ell\in I_j,\\
    z_\ell,&\ell\notin I_1\cup\cdots\cup I_m.
  \end{cases}
\]
The sets $I_1,\ldots,I_m$ are the \emph{blocks} of the map, and $mk/n$ is its \emph{active proportion}.

The action of $G$ on $X$ has the \emph{dense block property} if for every $r, m \in \mathbb{N}$ and every $0 < \eta < 1$, there are $n, k \in \mathbb{N}$ such that $mk \geq (1 - \eta)n$, and every coloring $c : X^n \to [r]$ admits a $G$-block map $\phi : X^m \to X^n$ of block size $k$ satisfying \[c(\phi(x)) = c(\phi(y)) \text{ for all $x, y \in X^m$ such that $y_j \in G \cdot x_j$ for every $j \in [m]$.}\]

Said differently, the induced coloring $c\circ\phi$ depends only on the $G$-orbit of each input coordinate, while the active blocks cover at least a $1-\eta$ proportion of the ambient coordinates.
For instance, observe that if the action of $G$ on $X$ is transitive, then $c(\phi(x))$ is constant over all $x \in X^m$.

The following theorem of Pálvölgyi shows that solvability of the acting group suffices to guarantee the dense block property.


\begin{theorem}[Pálvölgyi’s Dense Block theorem {\cite[Theorem 11]{Palvolgyi26}}]\label{lem:palvolgyi}
    Every action of a finite solvable group on a finite set has the dense block property.
\end{theorem}


We apply this theorem to the diagonal action of the dihedral group on pairs of cycle vertices to obtain a distance-profile-canonical restriction while preserving the gap constraints.

\begin{lemma}[Canonicalization on a torus]\label{lem:make-nice}
    For all $0 < \varepsilon < 1/2$ and $R, n \in \mathbb{N}$, there exists $N \in \mathbb{N}$ with the following property.
    For every $\varepsilon$-gap coloring $c : (\mathbb{Z}_R)^N \times (\mathbb{Z}_R)^N \to \{0, 1\}$, there exist two maps $\iota_1, \iota_2 : (\mathbb{Z}_R)^n \to (\mathbb{Z}_R)^N$ such that the coloring $\widehat{c} : (\mathbb{Z}_{R})^n \times (\mathbb{Z}_{R})^n \to \{0, 1\}$ defined by $\widehat{c}(x, y) = c(\iota_1(x), \iota_2(y))$ is an $\varepsilon/2$-gap distance-profile-canonical coloring.  
\end{lemma}

Note that allowing different row and column embeddings lets us obtain a distance-profile-canonical restriction without first symmetrizing the original coloring.


\begin{proof}
    Consider the dihedral group $D_R$ (the group of symmetries of the $R$-gon). 
    Note that $D_R \cong C_R \rtimes C_2$ is a finite solvable group.
    By identifying $\mathbb{Z}_R$ with the $R$-gon, $D_R$ naturally acts on $\mathbb{Z}_R$.
    Consider the diagonal action of $D_R$ on $\mathbb{Z}_R \times \mathbb{Z}_R$, defined by $g \cdot (x, y) = (g \cdot x, g \cdot y)$.
    Let $N, k$ be given by \cref{lem:palvolgyi} for this action of $D_R$ on $\mathbb{Z}_R \times \mathbb{Z}_R$, for $r=2$, $n$ and $\eta = \varepsilon/2$.
    In particular, we have $nk \geq (1-\varepsilon/2)N$.
    
    For every $m \geq 1$, we identify the sets $(\mathbb{Z}_R)^m \times (\mathbb{Z}_R)^m$ and $(\mathbb{Z}_R \times \mathbb{Z}_R)^m$ via the bijection \[((x_1, \ldots, x_m), (y_1, \ldots, y_m)) \mapsto ((x_1, y_1), \ldots, (x_m, y_m)).\]
    Let $c : (\mathbb{Z}_R)^N \times (\mathbb{Z}_R)^N \to \{0, 1\}$ be an $\varepsilon$-gap coloring.
    We view $c$ as a $2$-coloring of $(\mathbb{Z}_R \times \mathbb{Z}_R)^N$.
    Applying \cref{lem:palvolgyi} to the diagonal action of $D_R$ on $\mathbb{Z}_R \times \mathbb{Z}_R$ for $r=2$, $n$ and $\eta = \varepsilon/2$, we obtain a $D_R$-block map $\phi : (\mathbb{Z}_R \times \mathbb{Z}_R)^n \to (\mathbb{Z}_R \times \mathbb{Z}_R)^N$ of block size $k$ such that for all $(x, y), (x', y') \in (\mathbb{Z}_R \times \mathbb{Z}_R)^n$, we have $c(\phi(x, y)) = c(\phi(x', y'))$ whenever $(x'_i, y'_i) \in D_{R} \cdot (x_i, y_i)$ for every $i \in [n]$.
    
    Let $I_1, \ldots, I_n \subseteq [N]$ be the corresponding blocks, recall that each has cardinality $k$, let $g_{t} \in D_R$ for $t \in I_1 \cup \ldots \cup I_n$ be the corresponding labels, and $z_{t} = (a_t, b_t) \in \mathbb{Z}_R \times \mathbb{Z}_R$ for $t$ outside of this union be the corresponding fixed letters.
    
    Observe that for $(a, b), (a', b') \in \mathbb{Z}_R \times \mathbb{Z}_R$, we have $(a', b') \in D_R \cdot (a, b)$ if and only if $d_{\mathbb{Z}_R}(a, b) = d_{\mathbb{Z}_R}(a', b')$.
    Indeed, after rotating $a$ to $a'$, a reflection, if necessary, matches the orientation of the displacement from $a$ to $b$ with that from $a'$ to $b'$.
    Thus, for all $x, y, x', y' \in (\mathbb{Z}_R)^n$, we have $c(\phi(x, y)) = c(\phi(x', y'))$ whenever $\delta(x, y) = \delta(x', y')$.
    
    Since the action of $D_R$ on $\mathbb{Z}_R \times \mathbb{Z}_R$ is diagonal, the map $\phi : (\mathbb{Z}_R \times \mathbb{Z}_R)^n \to (\mathbb{Z}_R \times \mathbb{Z}_R)^N$ factors into two maps $\iota_1, \iota_2 : (\mathbb{Z}_R)^n \to (\mathbb{Z}_R)^N$ so that $\phi(x, y) = (\iota_1(x), \iota_2(y))$ for all $x, y \in (\mathbb{Z}_R)^n$.
    
    Thus, for all $x, y, x', y' \in (\mathbb{Z}_R)^n$ such that $\delta(x, y) = \delta(x', y')$, we have \[\widehat{c}(x, y) = c(\iota_1(x), \iota_2(y)) = c(\phi(x, y)) = c(\phi(x', y')) = c(\iota_1(x'), \iota_2(y')) = \widehat{c}(x', y').\]
    This proves that $\widehat{c}$ is distance-profile-canonical.
    
    Finally, let $x, y \in (\mathbb{Z}_R)^n$.
    Set $\Delta = \sum_{t \notin \bigcup_{j \in [n]}I_j}d_{\mathbb{Z}_R}(a_t, b_t)$.
    Using that $D_R$ acts on $\mathbb{Z}_R$ by isometries, we have \begin{align*}
        d(\iota_1(x), \iota_2(y)) &= \sum_{t \in [N]}d_{\mathbb{Z}_R}(\iota_1(x)_t, \iota_2(y)_t) \\
            &= \sum_{j \in [n]}\sum_{t \in I_j} d_{\mathbb{Z}_R}(g_t \cdot x_j, g_t \cdot y_j) + \sum_{t \notin \bigcup_{j \in [n]}I_j}d_{\mathbb{Z}_R}(a_t, b_t) \\
            &= \sum_{j \in [n]}\sum_{t \in I_j} d_{\mathbb{Z}_R}(x_j, y_j) + \Delta \\
            &= k \sum_{j \in [n]}d_{\mathbb{Z}_R}(x_j, y_j) + \Delta \\
            &= k \cdot d(x, y) + \Delta.
    \end{align*} 
    We want to apply \cref{lem:transfer-gap-coloring} to $((\mathbb{Z}_R)^N, d_{(\mathbb{Z}_R)^N})$ and $((\mathbb{Z}_R)^n, d_{(\mathbb{Z}_R)^n})$ with the maps $\iota_1, \iota_2 : (\mathbb{Z}_R)^n \to (\mathbb{Z}_R)^N$.
    For this, we need to verify that $k \cdot \diam((\mathbb{Z}_R)^n) \geq (1 - \varepsilon/2) \cdot \diam((\mathbb{Z}_R)^N)$.
    Indeed, we have \[k \cdot \diam((\mathbb{Z}_R)^n) = kn \cdot \diam(\mathbb{Z}_R) \geq (1 - \varepsilon/2) \cdot N \cdot \diam(\mathbb{Z}_R) = (1 - \varepsilon/2) \cdot \diam((\mathbb{Z}_R)^N).\]
    Thus, $\widehat{c}$ is an $\varepsilon/2$-gap coloring by \cref{lem:transfer-gap-coloring}.
\end{proof}

\begin{thebibliography}{99}
\bibitem[]{Palvolgyi26} Dömötör Pálvölgyi. A nearcircumsphere-{R}amsey Theorem for Solvable Transitive Configurations. (2026)
\end{thebibliography}
\end{document}
